\documentclass{article}
\usepackage{graphicx} % Required for inserting images
\usepackage{amsmath}
\usepackage{amssymb}
\usepackage{amsthm}
%\usepackage{hyperref}
\usepackage{xcolor} % Required for defining custom colors
\usepackage{datetime} % Add this line to include the datetime package

\definecolor{darkblue}{rgb}{0.0, 0.0, 0.5} % Define a darker blue color

\usepackage[colorlinks=true, linkcolor=darkblue, urlcolor=darkblue, citecolor=darkblue]{hyperref} % Use the darker blue color

\title{Subgroup}
\author{Aaron Honculada}
\date{\monthname[\month] \the\year}

\begin{document}

\maketitle

A \textbf{subgroup} is a subset of a group that is itself a group under the same operation.
\section{Definition}
Let $G$ be a group with operation $*$ and identity element $e$. A subset $H \subseteq G$ is called 
a \textbf{subgroup} of $G$ if:
\begin{enumerate}
    \item \textbf{Closure}: For all $a, b \in H$, we have $a * b \in H$.
    \item \textbf{Identity}: The identity element $e$ of $G$ is in $H$.
    \item \textbf{Inverses}: For each $a \in H$, the inverse $a^{-1}$ is also in $H$.
\end{enumerate}
If these conditions hold, we write $H \leq G$, meaning that $H$ is a subgroup of $G$.

\section{Examples}
\subsection{Integers in Rational Numbers}
\begin{itemize}
    \item $(\mathbb{Q}, +)$ is a group under addition.
    \item $\mathbb{Z}$ is a subgroup of $\mathbb{Q}$ because:
    \begin{itemize}
        \item \textbf{Closure}: For all $a, b \in \mathbb{Z}$, we have $a + b \in \mathbb{Z}$.
        \item \textbf{Identity}: The identity element $0 \in \mathbb{Z}$.
        \item \textbf{Inverses}: For each $a \in \mathbb{Z}$, the inverse $-a \in \mathbb{Z}$.
        \item So $\mathbb{Z} \leq \mathbb{Q}$ under addition.
    \end{itemize}
\end{itemize}

\subsection{Rotations in Symmetry Group \texorpdfstring{$S_3$}{S3}}
\label{subsec:rotations_s3}
The \textbf{symmetric group} $S_3$ is the group of all permutations of 3 elements.
\[
    S_3 = \{e, (12), (13), (23), (123), (132)\}
\]
The set of \textbf{only cyclic rotations} forms a subgroup:
\[
    H = \{e, (123), (132)\}
\]

\end{document}